\chapter{Rewriting with a Trace}\label{ch:rewrite}

Everything the engine does to a term, it does by rewriting: a rule looks at one node, and either
does not apply or replaces the node by something else and says why. The rewriter applies rules
until none applies, records every firing, and hands the record to the notebook. This chapter is
that machinery: the terms, the rules, the strategy, and the trace. It is deliberately small; the
weight of the book is in what the rules \emph{claim} (Part~II) and in why the rewriting
\emph{stops} (Part~III).

\section{Terms}

The term language is deliberately tiny.

\begin{minted}{lean}
inductive Expr where
  | num    : Q → Expr                    -- an exact rational, with an "approximate" flag
  | var    : String → Expr
  | add    : List Expr → Expr            -- n-ary sums and products
  | mul    : List Expr → Expr
  | pow    : Expr → Expr → Expr
  | fn     : String → List Expr → Expr   -- sin, ln, diff, matrix commands, …
  | matrix : List (List Expr) → Expr
\end{minted}

There is no subtraction and no division: $a - b$ is \lc{add [a, mul [-1, b]]} and $a/b$ is
\lc{mul [a, pow b (-1)]}. The printer recovers the familiar forms on the way out, and the
rewriter never has to know about them. Numerals are core Lean's \lc{Rat}, wrapped in a structure
\lc{Q} that adds a flag saying whether the number came from a decimal literal (so that
\texttt{2.5} prints as \texttt{2.5} and not \texttt{5/2}); the flag is presentation and changes no
value, and a two-line theorem in \file{proofs/Proofs/Q.lean} says so.

Sums and products are lists, and their arguments are kept in a canonical order. The function
\lc{canon} sorts the arguments of a sum or a product (products containing a matrix are left
alone: matrix multiplication does not commute). Canonical order runs \emph{silently} inside the
rewriter rather than as a rule: it is applied to every node after its children are normalized,
it is never recorded as a step, and Chapter~\ref{ch:origins} will have to bridge over it when it
traces origins across steps.

\section{Rules}

A rule inspects one node and either declines or returns a result: the new node, an explanation
in Markdown-with-\LaTeX, optionally a nested derivation (row reduction records its row operations
this way), and optionally an error --- a rule may refuse the whole evaluation, as the dimension
check in a matrix product does.

\begin{minted}{lean}
structure RuleResult where
  result : Expr
  explanation : String
  sub : Option Derivation := none
  error : Option String := none
\end{minted}

The rules of the original \lc{simplify} carry one more thing: a proof that they make a measure go
down.

\begin{minted}{lean}
structure Weights where
  w : Expr → Nat
  pos : ∀ e, 1 ≤ w e
  head : ∀ e cs, w (withChildren e cs) = w e   -- a weight looks only at the node's head

structure Rule (W : Weights) where
  name : String
  silent : Bool := false
  apply : Expr → Option RuleResult
  decreasing : ∀ e r, apply e = some r → measure W r.result < measure W e
\end{minted}

\lc{measure W} is the weighted node count. Because a weight depends only on a node's head and is
at least one, two facts are proved once for every weighting: the measure is additive over the
children (replacing children by lighter children makes the node lighter), and it is invariant
under \lc{canon} (a permutation of the arguments weighs the same). Those two facts are what let
the rewriter below be a total, well-founded function rather than a loop with a counter.

\begin{notebox}
  The obligation \lc{decreasing} is discharged by each rule's author, in Lean, at the point of
  writing the rule. For the seven original simplification rules the discharge is a few lines of
  \lc{omega} each under weights $2$ for a numeral, $4$ for a sum or product, $5$ for a power, $8$
  for a variable, a function or a matrix. That was enough for \lc{simplify}. It was not enough for
  derivatives, and Part~III is the story of what replaced it.
\end{notebox}

\section{The strategy}

The rewriter is \emph{innermost}: it normalizes the children of a node first, then applies the
first rule that fires at the node, then normalizes the result again. Innermost is not the only
choice, but it is the one that makes the later termination argument possible
(Chapter~\ref{ch:order}): a rule only ever fires on a node whose children are already normal.

\begin{minted}{lean}
mutual
  def normAt (rules : List (Rule W)) (e : Expr) (path : Path) (acc : Array RawStep) :
      {r : Expr × Array RawStep // measure W r.1 ≤ measure W e} := …
  def normChildren … := …
end

def normalize (rules : List (Rule W)) (e : Expr) : TraceM Expr := do
  let ⟨(out, raw), _⟩ := normAt rules e [] #[]
  let (steps, _) := buildSteps e raw
  modify (· ++ steps)
  pure out
\end{minted}

The return type of \lc{normAt} is a subtype: the function returns its result \emph{together with a
proof} that the result is no heavier than the input. That proof is what the recursion is measured
by --- after a rule fires, the result is strictly lighter, so normalizing it again is a smaller
problem --- and Lean accepts the definition without \lc{partial}, without fuel, and without any
promise the code does not keep.

\section{The trace}

A step records the rule, the explanation, the \emph{path} from the root to the node that changed,
and the whole term before and after.

\begin{minted}{lean}
structure Step where
  rule : String
  explanation : String
  path : Path            -- child indices from the root
  before : Expr
  after : Expr
  sub : Option Derivation := none

structure Derivation where
  input : Expr
  steps : Array Step
  output : Expr
\end{minted}

Paths are how the notebook and the engine refer to the same node. When the engine prints a term as
\LaTeX\ for the notebook it wraps every subterm in an annotation carrying its path,
\verb|\htmlData{path=1.0}{…}|, and KaTeX turns that into a DOM attribute. A click on the rendered
$\cos x$ yields the path \texttt{1.0}; the notebook sends it back with the cell and the line it was
clicked on; the engine walks its own copy of the term to the same node. No text is ever parsed
back. Chapter~\ref{ch:origins} builds on this to answer ``where did this come from?''.

\begin{logbox}[Parity, then a differential test]
  The engine replaced a TypeScript reference implementation. The port was tested against the
  reference at the wire level: 147 cell sources in one session, comparing text, \LaTeX\ with paths,
  value trees, errors and bindings --- zero mismatches, after which the reference was deleted and
  its answers kept as a golden file replayed by \texttt{lake test}. The differential test found
  two things worth keeping: matrix products must not be sorted into canonical order, and the two
  ``parity'' rules $(ab)^n \to a^n b^n$ and $(b^m)^n \to b^{mn}$ with symbolic $m$ resist every
  additive measure. They wait for Chapter~\ref{ch:order}.
\end{logbox}

\section{Exercises}

\begin{exercisebox}[A rule with its obligation]
  Write the rule $x + 0 \to x$ as a \lc{Rule simpW}: the \lc{apply} function on a node
  \lc{.add es}, removing the zero numerals, and the \lc{decreasing} proof. Under which weightings
  does the obligation hold? (It holds for every \lc{Weights}, since a node was removed; say why
  the \lc{head} law is what makes that a one-line proof.)
\end{exercisebox}

\begin{exercisebox}[Silent steps]
  Canonical ordering is not recorded as a step. Give a derivation in which the term after step $n$
  and the term before step $n{+}1$ differ, and describe what the notebook would show if the
  reordering \emph{were} recorded. Chapter~\ref{ch:origins} has to deal with exactly this gap.
\end{exercisebox}

\begin{exercisebox}[Paths]
  The path of the $\cos x$ in $2x\cos x + x^2\sin x$ depends on canonical order. Compute it from
  the tree \lc{add [mul [2, x, cos x], mul [x^2, sin x]]}, then change the order of the two
  summands and compute it again. Why must the engine, not the notebook, be the one to assign
  paths?
\end{exercisebox}
